\sectioncentered*{Заключение}
\addcontentsline{toc}{section}{Заключение}

В рамках работы была разработана информационно-поисковая система, реализующая модуль
индексирования документов на основе векторной модели с TF-IDF-взвешиванием и поиском
по косинусной близости. Система развёртывается в локальной вычислительной сети
посредством Docker Compose и предназначена для работы с документами на английском и французском языках.

В ходе выполнения работы были решены следующие задачи:

\begin{itemize}
	\item спроектирована и реализована трёхуровневая архитектура: статический веб-интерфейс (HTML/CSS/JavaScript), серверная часть на Python~3.11 (Flask) с разделением на модули по ответственности и уровень данных на PostgreSQL~16;
	\item спроектирована и создана база данных из семи таблиц (документы, логи поиска, словарь, IDF-значения, учёт миграций, клиенты-наблюдатели, логи событий наблюдателя) с применением расширения pgvector для хранения TF-IDF-векторов размерности 5000 и выполнения векторного поиска на стороне СУБД; в таблице документов предусмотрена колонка \texttt{lang} для хранения метки языка;
	\item реализованы основные алгоритмы системы: предобработка текста (нормализация, фильтрация стоп-слов, стемминг Портера для английского, сохранение диакритики для французского), построение словаря и вычисление инверсной документной частоты, TF-IDF-векторизация с логарифмическим сглаживанием и евклидовой нормализацией, ранжирующий поиск по косинусной близости с отбрасыванием нулевых результатов;
	\item реализована стратегия загрузки документов из файлов форматов .txt, .md, .pdf, .docx, .html, .rtf, .csv и .log с автоматическим извлечением текста и переиндексацией корпуса;
	\item реализован файловый наблюдатель (watchdog): CLI-клиент мониторит директорию и автоматически реплицирует создание, изменение и удаление файлов на сервер через REST API, с персистентным state-файлом (.watch\_state.json) для синхронизации при перезапусках;
	\item реализована подсветка релевантных терминов: как в сниппетах результатов поиска, так и при просмотре полного текста документа, с фолбэком на ключевые термины документа при отсутствии совпадений;
	\item реализована подсистема автоматического реферирования документов: два извлечительных метода (Sentence Extraction с TF-IDF и позиционными коэффициентами, TextRank с графовым ранжированием), многоязычная поддержка (EN/FR), REST API и веб-интерфейс с выбором метода и количества предложений;
	\item реализована страница оценки качества поиска: метрики точности, полноты и F-меры (F1, F0.5, F2) вычисляются по реальным результатам поиска и представляются аналитически (таблицей) и графически;
	\item реализована подсистема машинного перевода: пословый перевод EN$\to$FR
	      по словарю 95\,454 пары (41\,020 уникальных слов, партиционирован
	      по частям речи в PostgreSQL, покрывающий и трёхбуквенный GIN-индексы),
	      трёхуровневый lookup (in-memory $\to$ точный SQL $\to$ лемматизация spaCy),
	      нечёткий поиск по опечаткам (pg\_trgm), наполнение из English WordNet
	      $\times$ French OMW (сверх целевых 30\,000 пар),
	      сопряжение глаголов, французская грамматика (артикли, род, отрицание),
	      нейросетевой перевод (Helsinki-NLP/opus-mt-en-fr),
	      синтаксический разбор (spaCy + NLTK RegexpParser),
	      управление словарём, частотный список слов с переводом;
	      REST API и веб-интерфейс с выбором режима (Dictionary / Neural MT);
	\item реализована подсистема офлайн-синтеза речи (вариант~8) для
	      английского и французского языков: два движка "--- нейросетевой
	      Piper (4 ONNX-голоса, $\approx$63~МБ) и формантный eSpeak-NG
	      (3 голоса), 7 голосов с автоопределением языка (реюз подсистемы
	      распознавания), параметры темп/громкость/питч, нормализация
	      литературного текста, REST API (\texttt{/api/tts/*}) и
	      веб-интерфейс с выбором голоса и плеером;
	\item реализована подсистема офлайн-распознавания речи и голосового
	      управления (вариант~8) для английского и французского языков:
	      движок Vosk (Kaldi, 2 модели), двухпроходное распознавание
	      (free + constrained vocabulary), 16 фиксированных операций
	      с редактируемыми двуязычными триггер-фразами
	      (\texttt{/api/stt/commands}), реакция через существующие
	      сервисы (поиск, реферирование, перевод, детектор языка)
	      и навигацию, двуязычные уведомления toast + голос (реюз TTS),
	      запись с микрофона в браузере (16~кГц WAV), автотесты 25/25;
	\item обеспечено воспроизводимое развёртывание средствами Docker Compose с автоматическим применением миграций и инициализацией корпуса документов;
	\item реализован файловый наблюдатель (CLI-клиент на базе watchdog) для автоматической репликации изменений файловой системы на сервер с персистентным state-файлом.
\end{itemize}

Тестирование системы подтвердило работоспособность всех реализованных сценариев:
инициализации и миграций базы данных, поиска по запросам различной сложности, загрузки
документов в нескольких форматах, управления документами и оценки качества поиска.
В ходе тестирования выявлены и устранены дефекты, связанные с несоответствием контракта
клиент-серверного обмена при оценке метрик, обрезкой сниппетов на уровне SQL-запроса и
дублированием корпуса при повторных стартах.

Оценка качества на тестовой коллекции из 12 англоязычных документов по 11 тестовым
запросам показала среднюю полноту 1,0 и среднюю точность 0,302 при средней F1-мере
0,404. Анализ результатов показал, что при малом тематически однородном корпусе
качество поиска ограничивает точность: по общим запросам в топ-10 попадает много
нерелевантных документов с ненулевой косинусной близостью, тогда как по специфичным
запросам с редкими терминами система достигает идеальной точности. Система также
корректно работает с французскоязычными документами: язык документа классифицируется
автоматически при загрузке, язык запроса определяется при поиске, предобработка
выбирается в зависимости от языка.

Перспективными направлениями улучшения работы системы информационного поиска являются:
расширение тестового корпуса, переход к вероятностной модели ранжирования Okapi BM25
и гибридному (векторно-лексическому) поиску, замена стемминга лемматизацией, учёт фраз
и биграмм, автоматизация оценки качества с расширенным набором метрик (P@k, nDCG, MAP),
устранение фиксированной размерности вектора, а также перевод серверной части на
производственный WSGI-сервер.

Поставленная цель достигнута: разработанная система работоспособна, соответствует
векторной модели поиска и может использоваться как основа для дальнейших исследований
в области информационного поиска.